\section{Bölüm sonu çözümlü problemler}

\begin{enumerate}
  \item \textbf{Bağlı sıralar.} $2,4,4,7,9$ değerlerine sıra verin.

  \textbf{Çözüm:} Sıralar $1,2{,}5,2{,}5,4,5$'tir. Eşit iki 4, ikinci ve üçüncü sıraların ortalamasını alır.

  \item \textbf{Mann--Whitney etkisi.} $n_1=10$, $n_2=12$ ve birinci grup yönündeki $U=90$ ise üstünlük olasılığı ve sıra-biserial etkiyi bulun.

  \textbf{Çözüm:} $\hat\theta=U/(n_1n_2)=90/120=0{,}75$'tir. $r_{rb}=2\hat\theta-1=0{,}50$ olur. Birinci gruptan seçilen rastgele bir gözlemin ikinci gruptakinden yüksek olma eğilimi yüzde 75'tir; bağların yarısı hesaba katılır.

  \item \textbf{Kruskal--Wallis etkisi.} $H=12$, $k=3$ ve $N=36$ için epsilon-kareyi hesaplayın.

  \textbf{Çözüm:} $\varepsilon^2=(H-k+1)/(N-k)=(12-3+1)/(36-3)=10/33\approx0{,}303$'tür. Genel sonuç farkın hangi gruplar arasında olduğunu göstermez; düzeltilmiş ikili karşılaştırmalar gerekir.
\end{enumerate}